\chapter*{Akronieme}
\begin{acronym}[XXXX]
    \acro{API}{\textit{Application Programming Interface}}
    \acro{DB}{\textit{Database}}
    \acro{UI}{\textit{User Interface}}
    \acro{PWA}{\textit{Progressive Web App}}
    \acro{POPIA}{wetgewing op die beskerming van persoonlike inligting}
    \acro{LUKS}{\textit{Linux Unified Key Setup}}
    \acro{HTTPS}{\textit{Hypertext Transfer Protocol Secure}}
    \acro{CRC}{\textit{Class Responsibility Collaborator}}
    \acro{ERD}{\textit{Entity-Relationship Diagram}}
    \acro{SQL}{\textit{Structured Query Language}}
\end{acronym}

\section*{Skema vir vereiste nommering}

Die vereiste normering gebruik die eerste 4 karakters om die soort vereiste aan te dui en die volgende 3 karakters om die sub kategorie aan te dui.

\vspace{0.5em}
\begin{tabular}{@{}l@{\hspace{1cm}}l@{}}
    \textbf{FREQ-AAA-111} & Funksionele vereiste \\[0.3em]
    \textbf{NREQ-AAA-111} & Nie funksionele vereiste \\
\end{tabular}
\vspace{0.5em}

Waar AAA enige drie karakters is wat die kategorie aandui soos hieronder, en 111 'n unieke nommer is.

\vspace{0.5em}
\begin{tabular}{@{}l@{\hspace{1cm}}l@{}}
\textbf{USR} & Gebruikerbestuur \\[0.3em]
\textbf{MOB} & Mobiele toepassing \\[0.3em]
\textbf{WEB} & Web tuiste \\[0.3em]
\textbf{PRS} & Prestasie \\[0.3em]
\textbf{SCL} & Skaalbaarheid \\[0.3em]
\textbf{REL} & Betroubaarheid \\[0.3em]
\textbf{AVA} & Beskikbaarheid \\[0.3em]
\textbf{SEC} & Sekuriteit \\[0.3em]
\textbf{DAT} & Databestuur \\[0.3em]
\textbf{UIX} & Gebruikerskoppelvlak en -ervaring \\[0.3em]
\textbf{MNT} & Onderhoubaarheid \\[0.3em]
\textbf{TST} & Toets \\
\end{tabular}
\vspace{0.5em}
